
\medskip
\noindent
\underline{2. Fixed and toric parts of $T_\ell(A_K)$.  Criteria for
good reduction.  Orthogonality theorem for $\ell\ne p$.}

\medskip
\noindent
\SGARefTarget{sga7.sga7i.eix.n2-1}{2.1}. Like every commutative algebraic group over $k$, $A_o$ has a
largest subtorus, whose formation commutes with every extension of the
base field [SGA~3 XVII~7.2.1, XII~1.12]:
\begin{equation}\tag{\SGARefTarget{sga7.sga7i.eix.n2-1-1}{2.1.1}}
 T_o\subset A_o.
\end{equation}
When $k$ is perfect, $T_o$ is characterized by the condition that
$A_o^o/T_o$ be an extension of an abelian scheme $B_o$ over $k$ by a
smooth connected unipotent group [25].  The latter necessarily has a
composition series all of whose quotients are isomorphic to
$\mathbb G_{ak}$ (SGA~3 XVII~4.1.3), and one has an exact sequence
\begin{equation}\tag{\SGARefTarget{sga7.sga7i.eix.n2-1-2}{2.1.2}}
 0\longrightarrow U_o\longrightarrow A_o^o/T_o
 \longrightarrow B_o\longrightarrow0.
\end{equation}
In any case, whether or not $k$ is perfect, \underline{if} $A_o^o$ is
an extension of an abelian scheme $B_o$ by a smooth connected affine
scheme $L_o$,
\begin{equation}\tag{\SGARefTarget{sga7.sga7i.eix.n2-1-3}{2.1.3}}
 0\longrightarrow L_o\longrightarrow A_o^o
 \longrightarrow B_o\longrightarrow0,
\end{equation}
this extension structure is known to be unique [25].  Moreover,
$T_o\subset L_o$, and $T_o$ is the largest subtorus of $L_o$.  Thus
one has the exact sequence (\SGARefLink{sga7.sga7i.eix.n2-1-2}{2.1.2}), with
\begin{equation}\tag{\SGARefTarget{sga7.sga7i.eix.n2-1-4}{2.1.4}}
 U_o=L_o/T_o.
\end{equation}

\medskip
\noindent
\SGARefTarget{sga7.sga7i.eix.n2-1-5}{2.1.5}. Let us note that $A_o^o$ has \underline{unipotent rank zero}
(that is, after passage to an algebraic closure of $k$, it has no
subgroup isomorphic to $\mathbb G_a$) if and only if $A_o^o/T_o$ is an
abelian scheme $B_o$.  In that case the preceding hypotheses apply
with $U_o=0$.

Let $\ell$ be a prime different from the characteristic of $k$, so that
$A_o^o$ is $\ell$-divisible.  The inclusion (\SGARefLink{sga7.sga7i.eix.n2-1-1}{2.1.1}) then gives an exact
sequence of pro-group schemes:
\begin{equation}\tag{\SGARefTarget{sga7.sga7i.eix.n2-1-6}{2.1.6}}
\begin{tikzcd}[row sep=large,column sep=large]
0 \arrow[r] & T_\ell(T_o) \arrow[r]
 & T_\ell(A_o^o) \arrow[r] \arrow[d,"\wr"]
 & T_\ell(A_o^o/T_o) \arrow[r] \arrow[d,"\wr"] & 0\\
 & & T_\ell(A_o) & T_\ell(A_o/T_o) & {,}
\end{tikzcd}
\end{equation}
whose terms may also be interpreted as twisted constant free
$\ell$-adic sheaves.  When an exact sequence (\SGARefLink{sga7.sga7i.eix.n2-1-2}{2.1.2}) is available---for
example, when $k$ is perfect or when $A_o^o$ has unipotent rank
zero---one has
\begin{equation}\tag{\SGARefTarget{sga7.sga7i.eix.n2-1-7}{2.1.7}}
 T_\ell(A_o^o/T_o)\xrightarrow{\ \sim\ }T_\ell(B_o),
\end{equation}
so that the exact sequence (\SGARefLink{sga7.sga7i.eix.n2-1-6}{2.1.6}) becomes
\begin{equation}\tag{\SGARefTarget{sga7.sga7i.eix.n2-1-8}{2.1.8}}
 0\longrightarrow T_\ell(T_o)\longrightarrow T_\ell(A_o^o)
 \longrightarrow T_\ell(B_o)\longrightarrow0.
\end{equation}
Put
\begin{equation}\tag{\SGARefTarget{sga7.sga7i.eix.n2-1-9}{2.1.9}}
 n=\dim A_K=\dim A_o,\qquad
 \mu=\dim T_o,\qquad
 \lambda=\dim U_o,\qquad
 \alpha=\dim B_o.
\end{equation}
Then
\begin{equation}\tag{\SGARefTarget{sga7.sga7i.eix.n2-1-10}{2.1.10}}
 n=\mu+\lambda+\alpha,
\end{equation}
and
\begin{equation}\tag{\SGARefTarget{sga7.sga7i.eix.n2-1-11}{2.1.11}}
 \left\{
 \begin{aligned}
 \operatorname{rank}T_\ell(T_o)&=\mu,&
 \operatorname{rank}T_\ell(A_o^o/T_o)
 &=\operatorname{rank}T_\ell(B_o)=2\alpha,\\
 \operatorname{rank}T_\ell(A_o^o)&=\mu+2\alpha.
 \end{aligned}
 \right.
\end{equation}
The second relation is a well-known theorem of Weil [5, 15].

\medskip
\noindent
\SGARefTarget{sga7.sga7i.eix.n2-1-12}{2.1.12}. The integers $n,\mu,\lambda,\alpha$ are plainly defined, after
passing to an algebraic closure of $k$, even without assuming the
existence of an exact sequence (\SGARefLink{sga7.sga7i.eix.n2-1-2}{2.1.2}).  They are respectively called
the dimension of $A_o$, the \underline{reductive rank of $A_o$}
(equal to the dimension of $T_o$), the \underline{unipotent rank of
$A_o$}, and the \underline{abelian rank of $A_o$}.  They are still
related by (\SGARefLink{sga7.sga7i.eix.n2-1-10}{2.1.10}) and (\SGARefLink{sga7.sga7i.eix.n2-1-11}{2.1.11}), with the term
$\operatorname{rank}T_\ell(B_o)$ omitted from the latter.

\medskip
\noindent
\SGARefTarget{sga7.sga7i.eix.n2-1-13}{2.1.13}. Suppose that $k$ has characteristic $p>0$ and that $A_o^o$ is
$p$-divisible, or equivalently that it has unipotent rank zero
(cf. \SGARefLink{sga7.sga7i.eix.n2-1-5}{2.1.5}).  One should then also consider the pro-group $T_p(A_o^o)$,
which is a \underline{Barsotti--Tate group} (or a
\underline{$p$-divisible group}, in the terminology of [33]).  The
exact sequence (\SGARefLink{sga7.sga7i.eix.n2-1-8}{2.1.8}) remains valid for $\ell=p$ and is then an exact
sequence of Barsotti--Tate groups.  Here $\lambda=0$, and hence
\begin{equation}\tag{\SGARefTarget{sga7.sga7i.eix.n2-1-13-1}{2.1.13.1}}
 n=\mu+\alpha.
\end{equation}
The formulas (\SGARefLink{sga7.sga7i.eix.n2-1-11}{2.1.11}) also remain valid, provided that rank is
interpreted as the rank (or ``height'') of the Barsotti--Tate groups in
question.

\medskip
\noindent
\SGARefTarget{sga7.sga7i.eix.n2-1-14}{2.1.14}. For ${A'}_o^o$ we likewise introduce the invariants
\[
 T'_o\subset {A'}_o^o,\qquad U'_o,\qquad B'_o,
\]
where $T'_o$ is the toric part of ${A'}_o^o$, $B'_o$ its abelian part
(defined in particular when $k$ is perfect or when ${A'}_o^o$ has
unipotent rank zero), and $U'_o$ its unipotent part (defined, for
example, when $k$ is perfect, or trivially when ${A'}_o^o$ has
unipotent rank zero).  There is a corresponding exact sequence of the
form (\SGARefLink{sga7.sga7i.eix.n2-1-6}{2.1.6}), or (\SGARefLink{sga7.sga7i.eix.n2-1-8}{2.1.8}), which we shall not write again.  One could
introduce integers $\mu',\lambda',\alpha'$ for ${A'}_o^o$ as in
(\SGARefLink{sga7.sga7i.eix.n2-1-9}{2.1.9}), but we shall see below (\SGARefLink{sga7.sga7i.eix.n2-2-7}{2.2.7}) that they are the same as those
for $A_o^o$; this spares us from introducing new notation for them.
